% ===========================================================================
% ABSTRACT
% ===========================================================================

\begin{abstract}
Six machine learning surrogates for nonlinear Model Predictive Control
(MPC) -- LLNFM, LSTM, GP, PINN, TCN, and SW-MLP, evaluated as internal
predictors replacing an expensive physical plant model at every solver
step for wind turbine speed regulation in above-rated operation -- at
first appeared to hit a hard operational incompatibility with
SQP-based \texttt{nlmpc} for five of the six architectures. This paper
is a diagnostic study of that apparent incompatibility and a verified
remedy for it. A five-stage diagnostic traces the failure to
repeated \texttt{predict()} calls inside the solver --- not
architecture, model size, or MATLAB version. One-step prediction
accuracy and closed-loop tracking prove only loosely coupled. Independently,
\texttt{predict()}'s single-precision output was found to zero out
\texttt{fmincon}'s default finite-difference gradient estimate for
two of these architectures, corrupting solver feasibility regardless
of any real-time deadline. A verified manual forward
pass, computed in double precision, restored closed-loop viability for
all six surrogates. All
results hold under two explicit conditions: a bounded
thirty-iteration SQP solver budget, and a single software toolchain
(MATLAB Deep Learning Toolbox, tested on releases R2024a and R2026a) ---
no cross-toolchain replication was performed. On an Intel Core
i5-12450H (12th Gen, 16\,GB RAM) under R2024a, the headline result is a
$491\times$ mean per-step speedup for LSTM, the largest of the six
architectures tested (full range: $5.5\times$--$491\times$). We further identify a
Region~II/III torque-switching correction as a precondition for stable
closed-loop simulation, and, applying the same diagnostic rigor to our
own results, find that a previously reported "zero pitch activity"
result for a Mat\'ern-kernel GP uncertainty-penalized MPC cost is in
fact a solver-infeasibility artifact (a frozen actuator command, not
genuine regulation) rather than a controller achievement --- a finding
we report as a caution against under-instrumented closed-loop
comparisons, consistent with the diagnostic ethos of this paper.
These results reframe an apparent architectural limitation as a
diagnosable, remediable implementation cost, with practical guidance
for deploying ML surrogates in real-time MPC.

\textbf{Keywords:} Model Predictive Control, Wind Turbine Control,
Machine Learning Surrogates, Gaussian Process, Real-Time Feasibility
\end{abstract}
